\usepackage[spanish,es-noquoting]{babel}
\usepackage[utf8]{inputenc}
\usepackage[T1]{fontenc}
\usepackage{lmodern}
\usepackage{microtype}
\usepackage{amsmath,amssymb,amsfonts,amsthm,mathtools}
\usepackage{geometry}
\usepackage{xcolor}
\usepackage[most]{tcolorbox}
\usepackage{tikz}
\usetikzlibrary{arrows.meta,positioning,trees,calc,fit,shapes.geometric}
\usepackage{forest}
\usepackage{booktabs}
\usepackage{array}
\usepackage{enumitem}
\usepackage{hyperref}
\usepackage{cleveref}
\geometry{margin=2.6cm}

% Paleta suave
\definecolor{softBlue}{HTML}{EAF2FB}
\definecolor{softBlueLine}{HTML}{7EA6C8}
\definecolor{softGreen}{HTML}{ECF7EF}
\definecolor{softGreenLine}{HTML}{7FB58D}
\definecolor{softYellow}{HTML}{FFF8DC}
\definecolor{softYellowLine}{HTML}{D8B35A}
\definecolor{softRed}{HTML}{FBEAEA}
\definecolor{softRedLine}{HTML}{C97E7E}
\definecolor{softGray}{HTML}{F5F5F5}
\definecolor{softGrayLine}{HTML}{A0A0A0}

\newtheorem{definition}{Definición}[chapter]
\newtheorem{theorem}[definition]{Teorema}
\newtheorem{lemma}[definition]{Lema}
\newtheorem{proposition}[definition]{Proposición}
\newtheorem{corollary}[definition]{Corolario}
\newtheorem{invariant}[definition]{Invariante}

\theoremstyle{remark}
\newtheorem{remark}[definition]{Observación}
\newtheorem{example}[definition]{Ejemplo}

\newtcolorbox{intuitionbox}{colback=softBlue,colframe=softBlueLine,title=Intuición,breakable}
\newtcolorbox{formalbox}{colback=softGreen,colframe=softGreenLine,title=Formalización,breakable}
\newtcolorbox{warningbox}{colback=softRed,colframe=softRedLine,title=Cuidado,breakable}
\newtcolorbox{stepbox}[1]{colback=softYellow,colframe=softYellowLine,title={#1},breakable}
\newtcolorbox{summarybox}{colback=softGray,colframe=softGrayLine,title=Resumen,breakable}

\newcommand{\States}{\mathcal{S}}
\newcommand{\Moves}{\mathcal{M}}
\newcommand{\Legal}{\operatorname{Legal}}
\newcommand{\Succ}{\operatorname{Succ}}
\newcommand{\Term}{\operatorname{Terminal}}
\newcommand{\Player}{\operatorname{player}}
\newcommand{\Utility}{\operatorname{Utility}}
\newcommand{\Value}{\operatorname{Value}}
\newcommand{\Children}{\operatorname{Children}}
\newcommand{\MAX}{\mathrm{MAX}}
\newcommand{\MIN}{\mathrm{MIN}}
\newcommand{\true}{\mathrm{true}}
\newcommand{\false}{\mathrm{false}}

\newcommand{\blank}{\ensuremath{\cdot}}
\newcommand{\board}[9]{%
\begin{array}{|c|c|c|}\hline
#1 & #2 & #3 \\ \hline
#4 & #5 & #6 \\ \hline
#7 & #8 & #9 \\ \hline
\end{array}}

\forestset{
  game tree/.style={
    for tree={
      draw,
      rounded corners,
      align=center,
      edge={-Latex},
      parent anchor=south,
      child anchor=north,
      l sep=14mm,
      s sep=8mm,
      font=\small
    }
  }
}
